\chapter{行列式}
\section{习题}
\begin{enumerate}
  \item 判定以下 $9$ 阶排列的逆序数，从而判定它们的奇偶性：
  \begin{enumerate}
    \item $134782695$；
  \anssub
    \item $217986354$；
  \anssub
    \item $987654321$．
  \anssub
  \end{enumerate}
  \item 选择 $i$ 与 $k$，使
  \begin{enumerate}
    \item $1274i56k9$ 成偶排列；
  \anssub
    \item $1i25k4897$ 成奇排列．
  \anssub
  \end{enumerate}
  \item 写出把排列 $12435$ 变成排列 $25341$ 的那些对换．
  \anscompute
  \item 判定排列 $n(n-1)\cdots21$ 的逆序数，并讨论它的奇偶性．
  \anscompute
  \item 如果排列 $x_1x_2\cdots x_{n-1}x_n$ 的逆序数为 $k$，排列 $x_nx_{n-1}\cdots x_2x_1$ 的逆序数是多少？
  \ansgeneral
  \item 在 $6$ 阶行列式的展开式中，$a_{23}a_{31}a_{42}a_{56}a_{14}a_{65}$，$a_{32}a_{43}a_{14}a_{51}a_{66}a_{25}$ 这两项带有什么符号？
  \ansgeneral
  \item 写出 $4$ 阶行列式中所有带有负号并且包含因子 $a_{23}$ 的项．
  \anscompute
  \item 设
  \[
    A=
    \begin{vmatrix}
      a_{11}&a_{12}&a_{13}&a_{14}\\
      a_{21}&a_{22}&a_{23}&a_{24}\\
      a_{31}&a_{32}&a_{33}&a_{34}\\
      a_{41}&a_{42}&a_{43}&a_{44}
    \end{vmatrix}
  \]
  为 $4$ 阶行列式，$A_{ij}(1\leq i,j\leq 4)$ 为 $(i,j)$ 元素的代数余子式．写出下列算式为哪些行列式：
  \begin{enumerate}
    \item $\displaystyle\sum_{j_1j_2j_3j_4}(-1)^{\tau(j_1j_2j_3j_4)}a_{1j_1}b_{2j_2}a_{3j_3}a_{4j_4}$；% 原书此处印作 b，疑为 a 之误
  \anssub
    \item $\displaystyle\sum_{j_1j_2j_3j_4}(-1)^{\tau(j_1j_2j_3j_4)}a_{1j_1}a_{3j_2}a_{2j_3}a_{4j_4}$；
  \anssub
    \item $a_{31}A_{21}+a_{32}A_{22}+a_{33}A_{23}+a_{34}A_{24}$；
  \anssub
    \item $a_{13}A_{12}+a_{23}A_{22}+a_{33}A_{32}+a_{43}A_{42}$；
  \anssub
    \item $\displaystyle\sum_{j_1j_2j_3j_4}(-1)^{\tau(j_1j_2j_3j_4)}a_{j_1 1}b_{j_2 2}a_{j_3 3}a_{j_4 4}$，% 原书此处印作 b，疑为 a 之误
  \anssub
  \end{enumerate}
  其中 $\displaystyle\sum_{j_1j_2j_3j_4}$ 表示对所有 $4$ 阶排列的求和．
  \item 按定义计算下列行列式：
  \begin{enumerate}
    \item \[
    \begin{vmatrix}
      0&0&\cdots&0&1\\
      0&0&\cdots&2&0\\
      \vdots&\vdots&&\vdots&\vdots\\
      0&n-1&\cdots&0&0\\
      n&0&\cdots&0&0
    \end{vmatrix};
    \]
  \anssub
    \item \[
    \begin{vmatrix}
      0&1&0&\cdots&0\\
      0&0&2&\cdots&0\\
      \vdots&\vdots&\vdots&&\vdots\\
      0&0&0&\cdots&n-1\\
      n&0&0&\cdots&0
    \end{vmatrix};
    \]
  \anssub
    \item \[
    \begin{vmatrix}
      0&\cdots&0&1&0\\
      0&\cdots&2&0&0\\
      \vdots&&\vdots&\vdots&\vdots\\
      n-1&\cdots&0&0&0\\
      0&\cdots&0&0&n
    \end{vmatrix}.
    \]
  \anssub
  \end{enumerate}
  \item 由行列式定义证明：
  \[
    \begin{vmatrix}
      a_1&a_2&a_3&a_4&a_5\\
      b_1&b_2&b_3&b_4&b_5\\
      c_1&c_2&0&0&0\\
      d_1&d_2&0&0&0\\
      e_1&e_2&0&0&0
    \end{vmatrix}=0.
  \]
  \ansproof
  \item 由行列式定义计算
  \[
    f(x)=
    \begin{vmatrix}
      2x&x&1&2\\
      1&x&1&-1\\
      3&2&x&1\\
      1&1&1&x
    \end{vmatrix}
  \]
  中 $x^4$ 与 $x^3$ 的系数，并说明理由．
  \anscompute
  \item 由
  \[
    \begin{vmatrix}
      1&1&\cdots&1\\
      1&1&\cdots&1\\
      \vdots&\vdots&&\vdots\\
      1&1&\cdots&1
    \end{vmatrix}=0,
  \]
  证明奇、偶排列各半．
  \ansproof
  \item 设
  \[
    P(x)=
    \begin{vmatrix}
      1&x&x^2&\cdots&x^{n-1}\\
      1&a_1&a_1^2&\cdots&a_1^{n-1}\\
      \vdots&\vdots&\vdots&&\vdots\\
      1&a_{n-1}&a_{n-1}^2&\cdots&a_{n-1}^{n-1}
    \end{vmatrix},
  \]
  其中 $a_1,a_2,\cdots,a_{n-1}$ 是互不相同的数．
  \begin{enumerate}
    \item 由行列式定义，说明 $P(x)$ 是一个 $n-1$ 次多项式；
  \anssub
    \item 由行列式性质，求 $P(x)$ 的根．
  \anssub
  \end{enumerate}
  \item 计算下列行列式：
  \begin{enumerate}
    \item \[
    \begin{vmatrix}
      246&427&327\\
      1014&543&443\\
      -342&721&621
    \end{vmatrix};
    \]
  \anssub
    \item \[
    \begin{vmatrix}
      x&y&x+y\\
      y&x+y&x\\
      x+y&x&y
    \end{vmatrix};
    \]
  \anssub
    \item \[
    \begin{vmatrix}
      3&1&1&1\\
      1&3&1&1\\
      1&1&3&1\\
      1&1&1&3
    \end{vmatrix};
    \]
  \anssub
    \item \[
    \begin{vmatrix}
      1&2&3&4\\
      2&3&4&1\\
      3&4&1&2\\
      4&1&2&3
    \end{vmatrix};
    \]
  \anssub
    \item \[
    \begin{vmatrix}
      1+x&1&1&1\\
      1&1-x&1&1\\
      1&1&1+y&1\\
      1&1&1&1-y
    \end{vmatrix};
    \]
  \anssub
    \item \[
    \begin{vmatrix}
      a^2&(a+1)^2&(a+2)^2&(a+3)^2\\
      b^2&(b+1)^2&(b+2)^2&(b+3)^2\\
      c^2&(c+1)^2&(c+2)^2&(c+3)^2\\
      d^2&(d+1)^2&(d+2)^2&(d+3)^2
    \end{vmatrix}.
    \]
  \anssub
  \end{enumerate}
  \item 证明：
  \[
    \begin{vmatrix}
      b+c&c+a&a+b\\
      b_1+c_1&c_1+a_1&a_1+b_1\\
      b_2+c_2&c_2+a_2&a_2+b_2
    \end{vmatrix}
    =2
    \begin{vmatrix}
      a&b&c\\
      a_1&b_1&c_1\\
      a_2&b_2&c_2
    \end{vmatrix}.
  \]
  \ansproof
  \item 算出下列行列式的全部代数余子式：
  \begin{enumerate}
    \item \[
    \begin{vmatrix}
      1&2&1&4\\
      0&-1&2&1\\
      0&0&2&1\\
      0&0&0&3
    \end{vmatrix};
    \]
  \anssub
    \item \[
    \begin{vmatrix}
      1&-1&2\\
      3&2&1\\
      0&1&4
    \end{vmatrix}.
    \]
  \anssub
  \end{enumerate}
  \item 计算下列行列式：
  \begin{enumerate}
    \item \[
    \begin{vmatrix}
      1&1&1&1\\
      2&1&1&-3\\
      1&2&2&5\\
      4&3&2&1
    \end{vmatrix};
    \]
  \anssub
    \item \[
    \begin{vmatrix}
      1&\dfrac12&1&1\\[4pt]
      -\dfrac13&1&2&1\\[4pt]
      \dfrac13&1&-1&\dfrac12\\[4pt]
      -1&1&0&\dfrac12
    \end{vmatrix};
    \]
  \anssub
    \item \[
    \begin{vmatrix}
      0&1&2&-1&4\\
      2&0&1&2&1\\
      -1&3&5&1&2\\
      3&3&1&2&1\\
      2&1&0&3&5
    \end{vmatrix};
    \]
  \anssub
    \item \[
    \begin{vmatrix}
      1&\dfrac12&0&1&-1\\[4pt]
      2&0&-1&1&2\\
      3&2&1&\dfrac12&0\\[4pt]
      1&-1&0&1&2\\
      2&1&3&0&\dfrac12
    \end{vmatrix}.
    \]
  \anssub
  \end{enumerate}
  \item 计算下列 $n$ 阶行列式：
  \begin{enumerate}
    \item \[
    \begin{vmatrix}
      x&y&0&\cdots&0&0\\
      0&x&y&\cdots&0&0\\
      \vdots&\vdots&\vdots&&\vdots&\vdots\\
      0&0&0&\cdots&x&y\\
      y&0&0&\cdots&0&x
    \end{vmatrix};
    \]
  \anssub
    \item \[
    \begin{vmatrix}
      a_1-b_1&a_1-b_2&\cdots&a_1-b_n\\
      a_2-b_1&a_2-b_2&\cdots&a_2-b_n\\
      \vdots&\vdots&&\vdots\\
      a_n-b_1&a_n-b_2&\cdots&a_n-b_n
    \end{vmatrix};
    \]
  \anssub
    \item \[
    \begin{vmatrix}
      x_1-m&x_2&\cdots&x_n\\
      x_1&x_2-m&\cdots&x_n\\
      \vdots&\vdots&&\vdots\\
      x_1&x_2&\cdots&x_n-m
    \end{vmatrix};
    \]
  \anssub
    \item \[
    \begin{vmatrix}
      1&2&2&\cdots&2\\
      2&2&2&\cdots&2\\
      2&2&3&\cdots&2\\
      \vdots&\vdots&\vdots&&\vdots\\
      2&2&2&\cdots&n
    \end{vmatrix};
    \]
  \anssub
    \item \[
    \begin{vmatrix}
      1&2&3&\cdots&n-1&n\\
      1&-1&0&\cdots&0&0\\
      0&2&-2&\cdots&0&0\\
      \vdots&\vdots&\vdots&&\vdots&\vdots\\
      0&0&0&\cdots&n-1&1-n
    \end{vmatrix}.
    \]
  \anssub
  \end{enumerate}
  \item 证明：
  \begin{enumerate}
    \item \[
    \begin{vmatrix}
      a_0&1&1&\cdots&1\\
      1&a_1&0&\cdots&0\\
      1&0&a_2&\cdots&0\\
      \vdots&\vdots&\vdots&&\vdots\\
      1&0&0&\cdots&a_n
    \end{vmatrix}
    =a_1a_2\cdots a_n\left(a_0-\sum_{i=1}^n\frac{1}{a_i}\right);
    \]
  \anssub
    \item \[
    \begin{vmatrix}
      x&0&0&\cdots&0&a_0\\
      -1&x&0&\cdots&0&a_1\\
      0&-1&x&\cdots&0&a_2\\
      \vdots&\vdots&\vdots&&\vdots&\vdots\\
      0&0&0&\cdots&x&a_{n-2}\\
      0&0&0&\cdots&-1&x+a_{n-1}
    \end{vmatrix}
    =x^n+a_{n-1}x^{n-1}+\cdots+a_1x+a_0;
    \]
  \anssub
    \item \[
    \begin{vmatrix}
      \alpha+\beta&\alpha\beta&0&\cdots&0&0\\
      1&\alpha+\beta&\alpha\beta&\cdots&0&0\\
      0&1&\alpha+\beta&\cdots&0&0\\
      \vdots&\vdots&\vdots&&\vdots&\vdots\\
      0&0&0&\cdots&1&\alpha+\beta
    \end{vmatrix}
    =\frac{\alpha^{n+1}-\beta^{n+1}}{\alpha-\beta};
    \]
  \anssub
    \item \[
    \begin{vmatrix}
      \cos\alpha&1&0&\cdots&0&0\\
      1&2\cos\alpha&1&\cdots&0&0\\
      0&1&2\cos\alpha&\cdots&0&0\\
      \vdots&\vdots&\vdots&&\vdots&\vdots\\
      0&0&0&\cdots&1&2\cos\alpha
    \end{vmatrix}
    =\cos n\alpha;
    \]
  \anssub
    \item \[
    \begin{vmatrix}
      1+a_1&1&1&\cdots&1&1\\
      1&1+a_2&1&\cdots&1&1\\
      1&1&1+a_3&\cdots&1&1\\
      \vdots&\vdots&\vdots&&\vdots&\vdots\\
      1&1&1&\cdots&1&1+a_n
    \end{vmatrix}
    =a_1a_2\cdots a_n\left(1+\sum_{i=1}^n\frac{1}{a_i}\right).
    \]
  \anssub
  \end{enumerate}
  \item 已知
  \[
    \begin{vmatrix}
      x&y&z&x+y+z\\
      3&0&2&0\\
      1&1&1&1\\
      2&2&2&1
    \end{vmatrix}=\frac12,
  \]
  计算行列式
  \[
    \begin{vmatrix}
      x-y&y&z-x&x+y+z\\
      3&0&-1&0\\
      y-x&2-y&x-z&2-x-y-z\\
      3&2&-1&1
    \end{vmatrix}.
  \]
  \anscompute
  \item 在 $2n$ 阶行列式 $D_{2n}$ 中，主对角线上元素都是 $a$，次对角线（从右上角到左下角这条对角线）上元素都是 $b$，其余元素都是 $0$．计算行列式 $D_{2n}$．
  \anscompute
  \item 用克拉默法则解下列线性方程组：
  \begin{enumerate}
    \item \[
    \begin{cases}
      2x_1-x_2+3x_3+2x_4=6,\\
      3x_1-3x_2+3x_3+2x_4=5,\\
      3x_1-x_2-x_3+2x_4=3,\\
      3x_1-x_2+3x_3-x_4=4;
    \end{cases}
    \]
  \anssub
    \item \[
    \begin{cases}
      x_1+2x_2+3x_3-2x_4=6,\\
      2x_1-x_2-2x_3-3x_4=8,\\
      3x_1+2x_2-x_3+2x_4=4,\\
      2x_1-3x_2+2x_3+x_4=-8;
    \end{cases}
    \]
  \anssub
    \item \[
    \begin{cases}
      x_1+2x_2-2x_3+4x_4-x_5=-1,\\
      2x_1-x_2+3x_3-4x_4+2x_5=8,\\
      3x_1+x_2-x_3+2x_4-x_5=3,\\
      4x_1+3x_2+4x_3+2x_4+2x_5=-2,\\
      x_1-x_2-x_3+2x_4-3x_5=-3;
    \end{cases}
    \]
  \anssub
    \item \[
    \begin{cases}
      5x_1+6x_2=1,\\
      x_1+5x_2+6x_3=0,\\
      x_2+5x_3+6x_4=0,\\
      x_3+5x_4+6x_5=0,\\
      x_4+5x_5=1.
    \end{cases}
    \]
  \anssub
  \end{enumerate}
  \item 设 $a_1,a_2,\cdots,a_n$ 是数域 $P$ 中互不相同的数，$b_1,b_2,\cdots,b_n$ 是数域 $P$ 中任一组给定的数．用克拉默法则证明：存在唯一的数域 $P$ 上的多项式 $f(x)=c_0x^{n-1}+c_1x^{n-2}+\cdots+c_{n-1}$，使
  \[
    f(a_i)=b_i,\quad i=1,2,\cdots,n.
  \]
  \ansproof
  \item 设水银密度 $h$ 与温度 $t$ 的关系为
  \[
    h=a_0+a_1t+a_2t^2+a_3t^3.
  \]
  由实验测定得以下数据：
  \begin{center}
    \begin{tabular}{c|cccc}
      $t/^{\circ}\mathrm{C}$ & $0$ & $10$ & $20$ & $30$\\ \hline
      $h/(\mathrm{g}\cdot\mathrm{cm}^{-3})$ & $13.60$ & $13.57$ & $13.55$ & $13.52$
    \end{tabular}
  \end{center}
  求当 $t=15^{\circ}\mathrm{C},40^{\circ}\mathrm{C}$ 时水银密度（精确到小数点后两位）．
  \anscompute
\end{enumerate}
\section{补充题}
\begin{enumerate}
  \item 求
  \[
    \sum_{j_1j_2\cdots j_n}
    \begin{vmatrix}
      a_{1j_1}&a_{1j_2}&\cdots&a_{1j_n}\\
      a_{2j_1}&a_{2j_2}&\cdots&a_{2j_n}\\
      \vdots&\vdots&&\vdots\\
      a_{nj_1}&a_{nj_2}&\cdots&a_{nj_n}
    \end{vmatrix},
  \]
  这里 $\displaystyle\sum_{j_1j_2\cdots j_n}$ 是对所有 $n$ 阶排列求和．
  \anscompute
  \item 证明：
  \[
    \frac{\mathrm{d}}{\mathrm{d}t}
    \begin{vmatrix}
      a_{11}(t)&a_{12}(t)&\cdots&a_{1n}(t)\\
      a_{21}(t)&a_{22}(t)&\cdots&a_{2n}(t)\\
      \vdots&\vdots&&\vdots\\
      a_{n1}(t)&a_{n2}(t)&\cdots&a_{nn}(t)
    \end{vmatrix}
    =
    \sum_{j=1}^n
    \begin{vmatrix}
      a_{11}(t)&\cdots&\dfrac{\mathrm{d}}{\mathrm{d}t}a_{1j}(t)&\cdots&a_{1n}(t)\\
      a_{21}(t)&\cdots&\dfrac{\mathrm{d}}{\mathrm{d}t}a_{2j}(t)&\cdots&a_{2n}(t)\\
      \vdots&&\vdots&&\vdots\\
      a_{n1}(t)&\cdots&\dfrac{\mathrm{d}}{\mathrm{d}t}a_{nj}(t)&\cdots&a_{nn}(t)
    \end{vmatrix}.
  \]
  \ansproof
  \item 证明：
  \begin{enumerate}
    \item \[
    \begin{vmatrix}
      a_{11}+x&a_{12}+x&\cdots&a_{1n}+x\\
      a_{21}+x&a_{22}+x&\cdots&a_{2n}+x\\
      \vdots&\vdots&&\vdots\\
      a_{n1}+x&a_{n2}+x&\cdots&a_{nn}+x
    \end{vmatrix}
    =
    \begin{vmatrix}
      a_{11}&a_{12}&\cdots&a_{1n}\\
      a_{21}&a_{22}&\cdots&a_{2n}\\
      \vdots&\vdots&&\vdots\\
      a_{n1}&a_{n2}&\cdots&a_{nn}
    \end{vmatrix}
    +x\sum_{i=1}^n\sum_{j=1}^n A_{ij},
    \]
    其中 $A_{ij}$ 是 $a_{ij}$ 的代数余子式；
  \anssub
    \item \[
    \sum_{i=1}^n\sum_{j=1}^n A_{ij}=
    \begin{vmatrix}
      a_{11}-a_{12}&a_{12}-a_{13}&\cdots&a_{1,n-1}-a_{1n}&1\\
      a_{21}-a_{22}&a_{22}-a_{23}&\cdots&a_{2,n-1}-a_{2n}&1\\
      \vdots&\vdots&&\vdots&\vdots\\
      a_{n1}-a_{n2}&a_{n2}-a_{n3}&\cdots&a_{n,n-1}-a_{nn}&1
    \end{vmatrix}.
    \]
  \anssub
  \end{enumerate}
  \item 计算下列 $n$ 阶行列式：
  \begin{enumerate}
    \item \[
    \begin{vmatrix}
      1&2&3&\cdots&n\\
      2&3&4&\cdots&1\\
      3&4&5&\cdots&2\\
      \vdots&\vdots&\vdots&&\vdots\\
      n&1&2&\cdots&n-1
    \end{vmatrix};
    \]
  \anssub
    \item \[
    \begin{vmatrix}
      \lambda&a&a&a&\cdots&a\\
      b&\alpha&\beta&\beta&\cdots&\beta\\
      b&\beta&\alpha&\beta&\cdots&\beta\\
      b&\beta&\beta&\alpha&\cdots&\beta\\
      \vdots&\vdots&\vdots&\vdots&&\vdots\\
      b&\beta&\beta&\beta&\cdots&\alpha
    \end{vmatrix};
    \]
  \anssub
    \item \[
    \begin{vmatrix}
      x&a&a&\cdots&a&a\\
      -a&x&a&\cdots&a&a\\
      -a&-a&x&\cdots&a&a\\
      \vdots&\vdots&\vdots&&\vdots&\vdots\\
      -a&-a&-a&\cdots&-a&x
    \end{vmatrix};
    \]
  \anssub
    \item \[
    \begin{vmatrix}
      x&y&y&\cdots&y&y\\
      z&x&y&\cdots&y&y\\
      z&z&x&\cdots&y&y\\
      \vdots&\vdots&\vdots&&\vdots&\vdots\\
      z&z&z&\cdots&z&x
    \end{vmatrix};
    \]
  \anssub
    \item \[
    \begin{vmatrix}
      1&1&\cdots&1\\
      x_1&x_2&\cdots&x_n\\
      x_1^2&x_2^2&\cdots&x_n^2\\
      \vdots&\vdots&&\vdots\\
      x_1^{n-2}&x_2^{n-2}&\cdots&x_n^{n-2}\\
      x_1^n&x_2^n&\cdots&x_n^n
    \end{vmatrix}.
    \]
  \anssub
  \end{enumerate}
  \item 计算 $f(x+1)-f(x)$，其中
  \[
    f(x)=
    \begin{vmatrix}
      1&0&0&0&\cdots&0&x\\
      1&2&0&0&\cdots&0&x^2\\
      1&3&3&0&\cdots&0&x^3\\
      \vdots&\vdots&\vdots&\vdots&&\vdots&\vdots\\
      1&n&C_n^2&C_n^3&\cdots&C_n^{n-1}&x^n\\
      1&n+1&C_{n+1}^2&C_{n+1}^3&\cdots&C_{n+1}^{n-1}&x^{n+1}
    \end{vmatrix}.
  \]
  \anscompute
  \item 图 2.1 表示一电路网络，每条线上标出的数字是电阻（单位：$\Omega$），$E$ 点接地，由 $X,Y,U,Z$ 点通入的电流皆为 $100\,\mathrm{A}$，求这四点的电位．（用基尔霍夫定律．）

  （图 2.1 说明：正方形四个顶点按顺时针依次为 $X,Y,Z,U$，中心为 $E$ 且接地；各线段所标电阻（单位：$\Omega$）为 $R_{XY}=\dfrac12$，$R_{YZ}=\dfrac13$，$R_{ZU}=\dfrac14$，$R_{UX}=1$，$R_{XE}=\dfrac16$，$R_{YE}=\dfrac17$，$R_{UE}=\dfrac15$，$R_{ZE}=\dfrac18$．）% 原书图 2.1 为电路图，此处以文字描述
  \anscompute
  \item 设 $x_0,x_1,\cdots,x_{n-1}$ 为 $n$ 个文字，计算行列式
  \[
    \begin{vmatrix}
      x_0&x_{n-1}&\cdots&x_1\\
      x_1&x_0&\cdots&x_2\\
      \vdots&\vdots&&\vdots\\
      x_{n-1}&x_{n-2}&\cdots&x_0
    \end{vmatrix}.
  \]
  \anscompute
  \item 判断行列式
  \[
    \begin{vmatrix}
      1&2&\cdots&n-2&n-1&n\\
      2^2&3^2&\cdots&(n-1)^2&n^2&(n+1)^2\\
      3^3&4^3&\cdots&n^3&(n+1)^3&(n+1)^3\\
      \vdots&\vdots&&\vdots&\vdots&\vdots\\
      n^n&(n+1)^n&\cdots&(n+1)^n&(n+1)^n&(n+1)^n
    \end{vmatrix}
  \]
  是否等于 $0$．
  \anscompute
\end{enumerate}
